\section*{第 \ref{ch:g9:chromosomes} 章——染色体}

\begin{solution}{exo:g9:chromosomes:1}
\omterm{def:g6:cell-unit-of-life:parts}{细胞核}中呈丝状的结构——其内容物盘绕，
以便运输。分裂时可见，染色良好；
分裂间期松散，不可见。
\end{solution}

\begin{solution}{exo:g9:chromosomes:2}
每个体\omterm{def:g5:first-look-at-cells:cell}{细胞}含 46 条，两两配对成 23 对；\omterm{def:g8:sexual-reproduction:gametes}{配子}含 23 条——
每对各取一条，位于\omterm{def:g8:sexual-reproduction:gametes}{配子}中。
\end{solution}

\begin{solution}{exo:g9:chromosomes:3}
$23 + 23 = 46$：\omterm{def:g5:flower-to-fruit:steps}{受精}使配对\omterm{def:g7:organs-at-work:recovery}{恢复}，
每对中一条来自母亲，一条来自父亲——
子代的每个体\omterm{def:g5:first-look-at-cells:cell}{细胞}中都有父母双方的\omterm{def:g9:chromosomes:chromosomes}{染色体}。
\end{solution}

\begin{solution}{exo:g9:chromosomes:4}
若\omterm{def:g8:sexual-reproduction:gametes}{配子}未减半，每代的\omterm{def:g9:chromosomes:chromosomes}{染色体}数都会翻倍：
46 条的\omterm{def:g8:sexual-reproduction:gametes}{配子}产生 92 条的子代，再产生 184 条的孙代。
减数分裂使\omterm{def:g5:biodiversity:species}{物种}的\omterm{def:g9:chromosomes:chromosomes}{染色体}数稳定在 46。
\end{solution}

\begin{solution}{exo:g9:chromosomes:5}
XX（女孩的\omterm{def:g9:chromosomes:chromosomes}{核型}）与 XY（男孩的\omterm{def:g9:chromosomes:chromosomes}{核型}）。父亲的情况是：
\omterm{def:g4:animal-reproduction:sexual}{卵细胞}都带 X，\omterm{def:g8:reproductive-systems:male}{精子}带 X 或 Y 各占一半——
到达\omterm{def:g4:animal-reproduction:sexual}{卵细胞}的\omterm{def:g8:reproductive-systems:male}{精子}决定性别。
\end{solution}

\begin{solution}{exo:g9:chromosomes:6}
这是一组可识别的\omterm{def:g9:heredity-and-traits:traits}{性状}，称为唐氏综合征。它证明
\omterm{def:g9:chromosomes:chromosomes}{染色体}携带\omterm{def:g9:heredity-and-traits:traits}{性状}信息：一条\omterm{def:g9:chromosomes:chromosomes}{染色体}的数目
改变，许多\omterm{def:g9:heredity-and-traits:traits}{性状}随之改变。
\end{solution}

\begin{solution}{exo:g9:chromosomes:7}
在减数分裂中，\omterm{def:g8:sexual-reproduction:gametes}{配子}从每对中获得哪一条
是逐对独立决定的——
每个亲本可产生 $2^{23}$ 种可能的半数\omterm{def:g9:chromosomes:chromosomes}{染色体}组，
\omterm{def:g5:flower-to-fruit:steps}{受精}时两种组合相乘。
\end{solution}

\begin{solution}{exo:g9:chromosomes:8}
决定\omterm{def:g9:heredity-and-traits:traits}{性状}的是\omterm{def:g9:heredity-and-traits:heredity}{遗传}信息，而非\omterm{def:g9:heredity-and-traits:traits}{性状}本身：\omterm{def:g9:chromosomes:chromosomes}{染色体}携带指令——
结构紧凑，位于\omterm{def:g5:first-look-at-cells:cell}{细胞}中。未显现的\omterm{def:g9:heredity-and-traits:traits}{性状}：两条同源\omterm{def:g9:chromosomes:chromosomes}{染色体}都存在于
每个\omterm{def:g5:first-look-at-cells:cell}{细胞}中，无论表现出何种\omterm{def:g9:heredity-and-traits:traits}{性状}。不同子代之间的组合各不相同，
因为减数分裂每次都是每对各取一条，
且每次都独立进行。
\end{solution}

\begin{solution}{exo:g9:chromosomes:9}
每次\omterm{def:g5:flower-to-fruit:steps}{受精}都是一次全新的抛掷：\omterm{def:g8:reproductive-systems:male}{精子}中 X 与 Y 各占一半，
每次尝试都是如此，过去的出生不会影响这场竞赛，
连生三个女孩，只是三次公平的抛掷——
第四次同样公平。
\end{solution}

\begin{solution}{exo:g9:chromosomes:10}
（a）女孩或女性的体\omterm{def:g5:first-look-at-cells:cell}{细胞}：完整的 46 条，XX。（b）携带 Y 的\omterm{def:g8:reproductive-systems:male}{精子}\omterm{def:g5:first-look-at-cells:cell}{细胞}：
是\omterm{def:g8:sexual-reproduction:gametes}{配子}，\omterm{def:g9:chromosomes:chromosomes}{染色体}数为一半；与\omterm{def:g4:animal-reproduction:sexual}{卵细胞}\omterm{def:g5:flower-to-fruit:steps}{受精}会发育成男孩。
（c）多了第 21 号\omterm{def:g9:chromosomes:chromosomes}{染色体}的体\omterm{def:g5:first-look-at-cells:cell}{细胞}：47 条——
唐氏综合征的\omterm{def:g9:chromosomes:chromosomes}{核型}，
XX，女孩的\omterm{def:g9:chromosomes:chromosomes}{核型}。
\end{solution}

\begin{solution}{exo:g9:chromosomes:11}
女王提供的每个\omterm{def:g4:animal-reproduction:sexual}{卵细胞}都带 X——
她们为第 23 对所提供的，别无他选。X 或 Y 的差别
只出在国王的\omterm{def:g8:sexual-reproduction:gametes}{配子}中：如果一定要把硬币的结果归于某人，
那么它出在国王的\omterm{def:g8:reproductive-systems:male}{精子}中——
无论如何，这次抛掷只是偶然，不能归咎于任何人。
\end{solution}

\begin{solution}{exo:g9:chromosomes:12}
问题在于：同一套\omterm{def:g9:chromosomes:chromosomes}{染色体}，\omterm{def:g5:first-look-at-cells:cell}{细胞}类型却各不相同——
若每个\omterm{def:g5:first-look-at-cells:cell}{细胞}都含有同样的 46 条\omterm{def:g9:heredity-and-traits:heredity}{遗传}指令，为什么\omterm{def:g8:nervous-communication:neuron}{神经元}不是\omterm{def:g1:the-five-senses:sense}{皮肤}\omterm{def:g5:first-look-at-cells:cell}{细胞}？
解决的方向是：\omterm{def:g5:first-look-at-cells:cell}{细胞}之间的差别在于
\emph{读取}哪些指令，而非拥有哪些——
\omterm{def:g9:chromosomes:chromosomes}{染色体}组好比一座图书馆，\omterm{def:g5:first-look-at-cells:cell}{细胞}分化则是一种借阅策略。
\omterm{def:g9:chromosomes:chromosomes}{染色体}上哪些单位可被读取，
是下一章的问题。
\end{solution}

\begin{solution}{pb:g9:chromosomes:1}
\textbf{1.} A：男性或男孩的体\omterm{def:g5:first-look-at-cells:cell}{细胞}，46 条，XY。B：
女性或女孩的体\omterm{def:g5:first-look-at-cells:cell}{细胞}，46 条，XX。C：\omterm{def:g8:sexual-reproduction:gametes}{配子}——
23 条，每对各取一条；其 X 使它成为\omterm{def:g4:animal-reproduction:sexual}{卵细胞}，
或带 X 的\omterm{def:g8:reproductive-systems:male}{精子}。

\textbf{2.} \omterm{def:g8:sexual-reproduction:gametes}{配子}。它的 23 条体现了减数分裂的特点：
由 46 条的\omterm{def:g5:first-look-at-cells:cell}{细胞}出发，每对中各取一条，
没有任何一对是完整保留的。

\textbf{3.} 47 条，多出的一条位于第 21 对，第 23 对为 XX：
这是患唐氏综合征女孩的\omterm{def:g9:chromosomes:chromosomes}{核型}，
一条\omterm{def:g9:chromosomes:chromosomes}{染色体}的数目出错，导致发育改变。

\textbf{4.} 来自父亲：其\omterm{def:g8:reproductive-systems:male}{精子}决定带 X 还是带 Y。
\omterm{def:g5:flower-to-fruit:steps}{受精}时，两种情况的概率各为二分之一。

\textbf{5.} 因为身体由一个\omterm{prop:g8:sexual-reproduction:core}{受精卵}发育而来，
通过复制性分裂：每次分裂都复制\omterm{def:g9:chromosomes:chromosomes}{染色体}组，
并将一整套拷贝分给每个子\omterm{def:g5:first-look-at-cells:cell}{细胞}——
从\omterm{def:g1:the-five-senses:sense}{皮肤}\omterm{def:g5:first-look-at-cells:cell}{细胞}到\omterm{def:g8:nervous-communication:neuron}{神经元}，都是 46 条忠实的拷贝。

\textbf{6.} \omterm{def:g8:sexual-reproduction:gametes}{配子}的减数分裂及其规则：
从 23 对中各取一条，形成完整的半数\omterm{def:g9:chromosomes:chromosomes}{染色体}组，
没有任何一对重复，也没有任何一对缺失。

\textbf{7.} 孩子从其中获得一条 X：
无论另一侧送来哪一种\omterm{def:g8:sexual-reproduction:gametes}{配子}，
C 对第 23 对的贡献已经固定——
携带 X 的\omterm{def:g8:sexual-reproduction:gametes}{配子}无\omterm{def:g6:exploring-our-environment:conditions}{条件}地将 X 写入孩子的第 23 对。

\textbf{8.} 23 对\omterm{def:g9:chromosomes:chromosomes}{染色体}中的每一对都独立地贡献其中任一条，
$2 \times 2 \times \dots$，共 23 次，
即 $2^{23}$ 种半数\omterm{def:g9:chromosomes:chromosomes}{染色体}组，数以百万计，全部来自同一位亲本。
\omterm{def:g5:flower-to-fruit:steps}{受精}时，双亲各提供这样一种组合，
两个组合空间相乘，
这种组合的算术解释了兄弟姐妹为何从不重复。

\textbf{9.} 没有人做错任何事：多出的 21 号\omterm{def:g9:chromosomes:chromosomes}{染色体}
源于一次减数分裂的意外——一对同源\omterm{def:g9:chromosomes:chromosomes}{染色体}未能分开，
这是机制上的偶然，而非
父母的行为或价值所致。第 47 条\omterm{def:g9:chromosomes:chromosomes}{染色体}改变了
孩子的发育，但这并不改变孩子完全是他们的孩子——
\omterm{def:g9:heredity-and-traits:heredity}{遗传}咨询的结论与
\omterm{def:g5:first-look-at-cells:cell}{细胞}\omterm{def:g9:heredity-and-traits:heredity}{遗传}学的结论一致。

\textbf{10.} 二分之一，与每次出生时一样：
\omterm{def:g8:reproductive-systems:male}{精子}群体在每次\omterm{def:g5:flower-to-fruit:steps}{受精}时都是一半 X、一半 Y，
过去的结果不会在其中留下任何痕迹。
这一连串结果看似有意义，但机制本身没有记忆。

\textbf{11.} 性别在\omterm{def:g5:flower-to-fruit:steps}{受精}时决定，
取决于哪个\omterm{def:g8:reproductive-systems:male}{精子}先到达——
发生在父方\omterm{def:g8:sexual-reproduction:gametes}{配子}之中、竞争的最后几毫米里。
饮食既不影响\omterm{def:g8:reproductive-systems:male}{精子}群体中一半对一半的产生，
也不影响竞争的结果；该说法与可阻断或造成偏倚的环节毫无关联——
\cref{met:g8:contraception:evaluate} 中的问题 1，
再次失败。

\textbf{12.} \omterm{def:g9:heredity-and-traits:heredity}{遗传}依托于\omterm{def:g6:cell-unit-of-life:parts}{细胞核}中可数的\omterm{def:g9:chromosomes:chromosomes}{染色体}：
减半进入\omterm{def:g8:sexual-reproduction:gametes}{配子}，\omterm{def:g5:flower-to-fruit:steps}{受精}时\omterm{def:g7:organs-at-work:recovery}{恢复}；
计数出错时改变\omterm{def:g9:heredity-and-traits:traits}{性状}，第 23 对决定性别，
四张照片，一个载体。
\end{solution}
